\title{Large Language Models Can Automatically Engineer Features for Few-Shot Tabular Learning}

\begin{abstract}
Large Language Models (LLMs) exhibit exceptional capacity for addressing complex and previously unseen reasoning tasks, presenting significant promise for tabular learning—a critical component in numerous practical applications. In this study, we introduce a new in-context learning approach, \textsf{FeatLLM}, which leverages LLMs as feature engineers to create a transformed input dataset tailored for superior tabular prediction. The features synthesized by \textsf{FeatLLM} serve as the input for a straightforward downstream machine learning method, like linear regression, thereby enabling effective few-shot learning and robust predictive performance. At inference, our \textsf{FeatLLM} method solely relies on this lightweight model and the derived features, omitting the necessity to query the LLM for each new sample. Furthermore, the approach is designed to require only API-level interaction with LLMs, circumventing issues related to prompt length constraints. Extensive evaluations across a variety of tabular datasets from different fields show that \textsf{FeatLLM} produces high-value rules and substantially outperforms prior techniques, such as TabLLM and STUNT, by an average margin of 10\%.
\end{abstract}

\section{Introduction}
Large Language Models (LLMs) have, in recent years, demonstrated remarkable proficiency in producing accurate responses for novel tasks, even when specific fine-tuning is not performed~\cite{creswell2022selection,imani2023mathprompter,taylor2022galactica}. These models, pretrained on vast amounts of textual data, encapsulate broad knowledge that frequently substitutes for specialized domain expertise~\cite{Genomes2021,singhal2023large,wilcox2003role}, enabling widespread automation in fields ranging from dialogue analytics~\cite{finch2023leveraging} to program repair~\cite{fan2023automated}. By introducing task definitions and limited examples within the prompt, LLMs are capable of understanding the specific task context, recognizing relevant features, and selectively focusing on important instances~\cite{brown2020language,zhao2023survey}. This insight underscores the suitability of LLMs as effective feature engineering agents, capable of insightful data analysis.

The exploitation of LLMs’ background knowledge and logical reasoning has recently been applied to tabular learning. For example, pretrained medical knowledge, such as recognizing increased susceptibility to diseases among older populations, has become beneficial for predictive tasks within healthcare. Contemporary methodologies articulate task and feature explanations in plain language, convert tabular data for LLM ingestion, and subsequently utilize LLMs for processing~\cite{dinh2022lift,wang2023anypredict}. Downstream strategies typically consist of either prompt-based in-context learning~\cite{nam2023semi} or efficient parameter fine-tuning~\cite{dinh2022lift,hegselmann2023tabllm}. Leveraging the inherent knowledge stored within LLMs has led to superior outcomes, especially in limited data scenarios, consistently outperforming traditional tabular learning algorithms~\cite{hegselmann2023tabllm}.

Existing approaches to tabular learning that utilize LLMs face several drawbacks. Firstly, in end-to-end prediction scenarios, they necessitate running an LLM inference for each individual data point, which is resource-intensive. Achieving high accuracy with most techniques further requires fine-tuning the LLMs, an approach that becomes impractical when the model’s training infrastructure or tuning capabilities are unavailable—a challenge that is exacerbated by the recent trend of high-performing LLMs being accessible solely via inference APIs with limited flexibility~\cite{anil2023palm,openai2023gpt4}. Additionally, the majority of these methods are ill-suited to handle prompts of considerable length~\cite{liu2023lost}; as the dimensionality of the tabular data increases and thus the serialized textual representation grows longer, either computational costs become prohibitive or performance deteriorates. Together, these issues create significant obstacles for the effective real-world adoption of LLM-based methods for tabular data analysis.

Our strategy represents a conceptual shift from using LLMs to directly predict target outputs. Instead, we repurpose LLMs primarily for feature engineering, capitalizing on their ability to encode vast prior knowledge. More concretely, we present a new in-context learning paradigm, called \textsf{FeatLLM}, focused on uncovering the underlying “criteria” that inform LLM predictions. For example, in a disease prediction context, our approach enables the LLM to extract and articulate explicit feature-based decision rules that indicate the presence or absence of a condition. These criteria are then translated into engineered features that substitute for the original attributes in downstream machine learning tasks. With the responsibility for generating meaningful features shifted to the LLM, subsequent learning stages can rely on simpler models, reducing the latency of inference processes. Utilizing the in-context capabilities of LLMs, we can obtain these features by crafting prompts with illustrative examples—sidestepping the need for further model training. Moreover, by incorporating ensemble learning concepts such as feature bagging~\cite{breiman2001random}, our method also addresses scalability concerns related to overly long prompts.

\textsf{FeatLLM} organizes the process of prompt formulation for the generation of new features into two primary sub-components: (1) comprehending the given problem and determining how the features are connected to the prediction targets; and (2) leveraging insights from this analysis, alongside a small number of example cases, to formulate rules that can effectively separate distinct classes. This systematic reasoning methodology guides LLMs to disregard features lacking informative value when formulating rules. These generated rules are then systematically implemented on the dataset (by interpreting them as executable program code), converting the data into binary indicators that reflect whether individual samples conform to each specific rule. Subsequently, a simple, low-complexity model is employed to infer class probabilities from these newly derived binary features. To further enhance robustness, an ensemble technique is adopted—new features are repeatedly constructed using a feature bagging scheme. By suitably tuning the quantity of features and the subset of samples in bagging, the limitations imposed by oversized prompts can be effectively alleviated. Since \textsf{FeatLLM} does not require any training, it incurs minimal inference costs, making it practical for diverse real-world deployment scenarios.

The performance of \textsf{FeatLLM} is empirically assessed across 13 distinct tabular datasets under settings with limited labeled samples, demonstrating that it consistently delivers reliable and competitive results. We provide evidence that LLMs are adept at leveraging both domain knowledge and patterns drawn from data to produce high-value features. Our approach surpasses the latest few-shot learning baselines in multiple experimental setups. The code is publicly available at the following anonymized GitHub repository:~\texttt{https://github.com/Sungwon-Han/FeatLLM}.

\section{Related Work}

\subsection{Few-Shot Learning with Tabular Data} 

Few-shot learning has made it possible to acquire representations that generalize well from data with limited supervision, significantly reducing annotation costs~\cite{chen2018closer}. While originally geared toward the image domain, the scope of few-shot learning has broadened to cover a variety of modalities such as text, audio, and more recently, tabular data~\cite{majumder2022few,nam2023stunt,schick2021exploiting}. Relying on very limited labeled samples, however, comes with the drawback of potential overfitting to coincidental patterns~\cite{bartlett2021deep}. To address this, contemporary research has explored leveraging auxiliary data or domain expertise to instill suitable inductive biases during training. Approaches in this direction include utilizing abundant unlabeled data—augmented in various ways—to support semi-supervised learning~\cite{ucar2021subtab,yoon2020vime}, as well as employing task-independent unsupervised pre-training to achieve more effective initializations~\cite{somepalli2021saint}. Common unsupervised learning methodologies for tabular data entail masking or intentionally corrupting portions of data and using reconstruction as a learning goal~\cite{arik2021tabnet,majmundar2022met}, or applying data transformations to facilitate contrastive learning frameworks~\cite{bahri2022scarf,chen2020simple}. Nevertheless, due to the inherent diversity of tabular data structures, finding universally robust augmentations is difficult, and such techniques generally yield limited gains in few-shot regimes~\cite{nam2023stunt}.

More recently, research has examined the strategy of training models on a diverse portfolio of tabular datasets—not necessarily related to the target application—and then deploying transfer learning for adaptation~\cite{zhang2023generative,levin2022transfer,zhu2023xtab}. An illustrative example is TransTab~\cite{wang2022transtab}, which adopts transformer architectures to develop semantic understanding between columns by exploiting column headers in addition to cell values. Through this cross-table and cross-column knowledge acquisition, these methods enable robust zero-shot or few-shot transfer to novel tabular tasks. Another instance, TabPFN~\cite{hollmann2022tabpfn}, constructs synthetic datasets by blending a mixture of distributions and utilizes them to pre-train a Bayesian neural network. In the inference stage, the model requires no further training: both training and test examples from the new task are directly input. The resultant prior derived from a diverse training set allows these models to outperform conventional tree-based algorithms and offers notable improvements when labeled data are scarce.

\subsection{Language-Interfaced Tabular Learning}

Incorporating large language models (LLMs)~\cite{anil2023palm,openai2023gpt4} presents an alternative and promising avenue for infusing prior knowledge into tabular learning. Trained on massive, varied text datasets, LLMs exhibit impressive generalization to a wide array of novel tasks, which has been demonstrated across multiple application domains~\cite{brown2020language}. Recent work has sought to capitalize on the embedded prior knowledge within LLMs by adapting them for tabular data challenges. Common practice involves converting tabular structures into serialized textual form to serve as prompt inputs for LLMs~\cite{dinh2022lift,hegselmann2023tabllm,wang2023anypredict}. By embedding details of the tabular data, explanatory information about features, and specific descriptions of the task within prompts, these systems enable LLMs to model task-feature dependencies during prediction. Advances in this area include either applying targeted, parameter-efficient fine-tuning techniques such as LoRA~\cite{hu2021lora} or IA3~\cite{liu2022few}, or alternatively, exploiting in-context learning by providing a handful of annotated examples within the prompt to facilitate learning without additional parameter updates~\cite{dinh2022lift,hegselmann2023tabllm,nam2023semi,slack2023tablet}.

Although the previously mentioned approaches have shown strong performance in enhancing few-shot tabular learning, their reliance on running inference directly through the LLM for each prediction creates limitations when deployed in industrial settings. This work moves away from the traditional paradigm where every sample must be processed end-to-end using an LLM, which treats the model as a black box. Instead, it instructs the LLM to derive explicit conditions or rule sets corresponding to each class. \textsf{FeatLLM} then translates these rules into programmatic code for feature construction, thereby enabling prediction without routing data through the LLM. In this setup, in-context learning is leveraged so that the model can extract relevant rules based on its prior knowledge as well as a small collection of labeled instances.

\section{Methods}
\textsf{FeatLLM} functions by identifying “rules,” i.e., defining conditions relevant to each output class, from the limited available samples, as illustrated in Figure~\ref{fig:main_model}. These conditions are first articulated through an LLM and subsequently converted into binary feature vectors for a given instance—each binary indicator reflecting whether its respective rule is met by the data point. The complete set of these generated binary features is then aggregated via a dot product with non-negative, learnable weights to compute class probabilities. The training phase updates these weights to quantify each rule’s significance regarding class likelihood (Section~\ref{Sec:training}). This entire methodology is conducted multiple times using bagging, and the results from each iteration are ultimately fused via ensemble techniques. Further elaboration on the problem statement and step-by-step details are provided below.

\vspace*{-0.12in}
\paragraph{Problem formulation.} Consider a tabular dataset $\mathcal{D} = \{(\mathbf{x}^i, \mathbf{y}^i)\}_{i=1}^{N}$ with $N$ labeled examples. Each entry $\mathbf{x}^i$ is a $d$-dimensional feature vector, and the dataset includes a set of natural language feature names $F = \{f_j\}_{j=1}^{d}$ and corresponding feature definitions presented independently. In this classification context, each label $\mathbf{y}^i$ belongs to a predefined class set. For $k$-shot learning scenarios, the training set is limited to $k$ randomly chosen labeled examples, where $k<N$.

\subsection{Prompt Design for Extracting Rules}
\label{Sec:prompt_design}
To enable the LLM to derive rules based on a logical reasoning process, our prompt construction intentionally directs the LLM to replicate the stepwise thinking that a domain expert would use in analyzing tabular data. The designed prompt consists of three primary segments, as illustrated in Fig.~\ref{fig:prompt_for_rule} and further detailed in Appendix~\ref{sec:example_rule}:

\vspace*{-0.12in}
\paragraph{Basic information description.} This segment communicates the vital background needed to approach the task (highlighted in orange in Fig.~\ref{fig:prompt_for_rule}). It covers a textual outline of the objective to be addressed, label context, feature explanations, and select illustrative examples. The task statement takes the form of a direct question (e.g., ``Does this patient's myocardial infarction complications data indicate chronic heart failure? Yes or no?"). Additional formulations can be found in Appendix~\ref{sec:task_desc}. Each feature is described by its value type, and may include further elaboration; for categorical variables, representative categories can be listed. To furnish exemplars, several training examples with their true labels are rendered as serial text, following established formats from prior works~\cite{dinh2022lift,hegselmann2023tabllm}:

\begin{align}
&\text{Serialize}(\mathbf{x}^i,\mathbf{y}^i, F) = \nonumber \\ 
&``\ f_1\ \text{is}\ \mathbf{x}^i_1.\ ... \ f_d\  \text{is}\  \mathbf{x}^i_d.\ \ \text{Answer:}\  \mathbf{y}^i\ ”
\end{align}

\vspace*{-0.12in}
\paragraph{Reasoning instruction.} Our objective is to improve the LLM's reasoning ability by furnishing targeted reasoning instructions (refer to the blue annotations in Fig.~\ref{fig:prompt_for_rule}). The reasoning instruction begins with an opening statement modeled on the chain-of-thought methodology~\cite{wei2022chain}. Subsequently, we segment the inference workflow into two stages. Initially, the LLM is prompted to deduce causal links or trends between the input features and the task specifications without consulting the demonstration examples, instead leveraging world knowledge or intuition to comprehensively retrieve its background knowledge concerning the problem. In the follow-up stage, the LLM integrates both the sample demonstrations and insights obtained from the first step to synthesize classification rules specific to each class. This sequential, two-phase reasoning helps the model avoid detecting coincidental associations in extraneous columns, thereby concentrating its attention on the most salient features. To circumvent issues such as insufficient or excessive rule extraction—which could either compromise accuracy through underfitting or inflate the prompt size—we enforce a fixed quota of rules to be extracted per class (specifically, 10)\footnote{A discussion regarding the effect of the number of rules appears in Section~\ref{sec:analysis}.}. Furthermore, to reduce type-related errors in rule formation, we instruct the LLM to consider the feature type as defined in the fundamental information section.

\vspace*{-0.12in}
\paragraph{Response instruction.} To support more straightforward rule parsing and downstream usage, we provide explicit instructions to the LLM on how to organize its generated responses (see yellow segments in Fig.~\ref{fig:prompt_for_rule}). These prompts delineate the expected response template for each predicted class (i.e., required response formatting) as well as the specific format that each rule should adhere to (i.e., rule formatting structure). This guidance empowers the LLM to select proper output formats based on the characteristics of each feature. In the absence of supplementary guidelines, the LLM retains the flexibility to connect several rules by employing logical connectors such as AND or OR. A sample of the resulting output can be found in Appendix~\ref{sec:outcome-rule}.

\subsection{Inferring Class Likelihood via Rules}
\label{Sec:training}

\paragraph{Parsing rules for feature generation.} To construct new binary features, we employ the rules obtained from the earlier step, generating features that correspond to each class. These binary features encode whether or not a sample meets the specified rules of that class (assigned a value of '0' or '1'). Such features can subsequently inform the prediction regarding a sample’s class membership. However, since the Large Language Model (LLM) produces rules formulated in natural language, an effective parsing approach is essential to enable automated feature extraction. Despite setting formatting constraints during the rule generation process to facilitate easier parsing, LLM outputs may still display variability and syntactic inconsistencies~\cite{kovavc2023large}. As an example, the LLM might substitute symbols like “$>$” or “$<$” with equivalent verbal descriptions such as “greater than” or “less than.” This inconsistency further complicates reliable parsing.

Rather than crafting complicated code solutions to manage noisy textual output, we capitalize on the strengths of the LLM itself to resolve these issues. The LLM’s capacity for semantic interpretation allows it to comprehend rule meaning even amidst syntactic divergence, and it can produce code snippets from textual instructions owing to its exposure to code datasets. To harness these abilities, our prompts to the LLM include the specific function name, explanations of input and output formats, and the set of inferred rules. We then provide these prompts to the LLM for code generation. Examples of such prompts, along with their results, are depicted in Figures~\ref{fig:prompt_for_parsing} and ~\ref{sec:example_parsing}-\ref{sec:example_parse_outcome} in the Appendix. The resultant code is executed via Python’s \texttt{exec()} command to facilitate data transformation.

\vspace*{-0.12in}
\paragraph{Inferring class likelihood.} After obtaining binary features derived from the rules per class, the class likelihood for any sample can be simply approximated by counting the number of applicable rules a sample fulfills for each class (i.e., summing the binary features corresponding to each class). Nevertheless, since the contribution of each rule or feature might differ, it becomes important to quantitatively estimate their relevance through learning from training data. To this end, we train a standard machine learning (ML) model to determine the relative weights of these binary features for every class. For instance, a linear model without a bias term could be employed. Let the binary feature vector produced for class $k$ by sample $\mathbf{x}^i$ be denoted $\mathbf{z}_k^i \in \{0, 1\}^R$, with $R$ representing the rule count per class, and $c$ denoting the total number of classes. The class probability vector $\mathbf{p}^i$ for each sample can then be derived as follows:

\begin{align}
\text{logit}_k^i &= \max(\mathbf{w}_k, 0) \cdot \mathbf{z}_k^i, \nonumber \\
\mathbf{p}^i &= \text{Softmax}({[\text{logit}_{1}^i, ..., \text{logit}_{c}^i]}). \label{eq:infer_prob}
\end{align}

In this context, $\mathbf{w}_k$ denotes the learnable parameter vectors within the linear model, which quantify each feature's relative contribution. By constraining these weights to the non-negative domain, we guarantee that LLM-derived features cannot negatively influence the prediction probability of any class during the training phase. Subsequently, the logit outputs are mapped to class probabilities via the softmax operation. The parameters $\mathbf{w}_k$ are estimated by minimizing the cross-entropy loss, utilizing a limited set of labeled examples.

\vspace*{-0.12in}
\paragraph{Ensembling with bagging.} To construct the ensemble, the described procedure is run multiple times, yielding several inference models. The ultimate class prediction is generated by averaging the class probabilities $\mathbf{p}^i$ produced by each set of weights $\mathbf{w}[t]$ from $T$ independent runs, following the protocol outlined in Eq.~\ref{eq:infer_prob}.

To foster diversity during the rule creation phase, a high temperature is chosen for LLM inference, which induces variability in generated rules. Furthermore, the sequence of few-shot examples is randomized; prior work has shown that example ordering can meaningfully affect the LLM's outputs~\cite{lu2022fantastically}\footnote{Detailed analysis regarding rule diversity is included in Appendix~\ref{sec:diversity}}. For additional diversity and stability in the aggregate predictions, we employ bagging—randomly subsampling subsets of instances or features for each ensemble member. This mechanism is particularly advantageous as the dataset or feature set grows in size. The ensemble paradigm conveys two key benefits: (1) errors from inaccurately generated rules are mitigated, as correct rules from other ensemble members predominate, yielding increased overall robustness. This advantage aligns with the principle of LLM self-consistency~\cite{wang2022self}, where rules that reappear across runs are more likely to be trustworthy. (2) Ensemble-based bagging also helps manage input length restrictions imposed by the LLM; when input tables become too extensive to fit within prompt limitations, subsampling preserves reliability without sacrificing performance. While any single model's generated rule set may omit some essential relationships, aggregating multiple weak learners ensures that missing information in one trial is compensated for by others, improving model robustness and coverage.

\section{Experiments}
We assess the effectiveness of \textsf{FeatLLM} on a variety of tabular datasets and perform comprehensive analyses to better understand the mechanisms underlying our framework. Owing to limitations in space, further experiments, such as tests on alternative LLMs and in-depth qualitative investigations, are included in the Appendix.

\subsection{Performance Evaluation}
\paragraph{Datasets.}
Our experimental setup makes use of 13 datasets focused on binary and multi-class classification problems: (1) Adult~\cite{asuncion2007uci}, predicting if an individual’s income surpasses \$50,000 per year; (2) Bank~\cite{moro2014data}, for anticipating if a customer will subscribe to a term deposit; (3) Blood~\cite{yeh2009knowledge}, determining the likelihood of a donor returning for subsequent donations; (4) Car~\cite{kadra2021well}, for classifying car quality; (5) Communities~\cite{misc_communities_and_crime_183}, for estimating crime rates in specified geographic regions; (6) Credit-g~\cite{kadra2021well}, which predicts the probability of an individual being a good or bad credit risk; (7) Diabetes\footnote{\texttt{kaggle.com/datasets/uciml/pima-indians-diabetes-database}}, for detecting diabetes occurrence in subjects; (8) Heart\footnote{\texttt{kaggle.com/datasets/fedesoriano/heart-failure-prediction}}, for identifying instances of coronary artery disease; (9) Myocardial~\cite{misc_myocardial_infarction_complications_579}, which assesses the likelihood of chronic heart failure in a given patient.

In addition, we incorporate newer and synthetic datasets, which are seldom discussed and do not overlap with any pre-training data for LLMs: (10) Cultivars, predicting high-yield potential in soybean cultivars\footnote{\texttt{archive.ics.uci.edu/dataset/913}}; (11) NHANES, for classifying individuals into age categories from available records\footnote{\texttt{archive.ics.uci.edu/dataset/887}}; (12) Sequence-type, distinguishing among four sequence classes: arithmetic, geometric, Fibonacci, or Collatz; and (13) Solution-mix, identifying if a solution composed of four mixtures surpasses 0.5 in concentration. Cultivars and NHANES are recently contributed datasets (post-September 2023) to the UCI Machine Learning Repository, confirming their exclusion from the LLM API’s (gpt-3.5-turbo-0613) pre-training datasets. Sequence-type and Solution-mix are entirely synthetic, ensuring zero exposure to the LLM before evaluation.

These datasets differ in both scale and complexity, with relevant statistics summarized in Table~\ref{Tab:basic-desc}. Every dataset supplies clear attribute names and descriptions. Notably, `Communities' and `Myocardial' possess more than 100 features, resulting in additional challenges for LLMs that operate under fixed context window constraints.

\vspace*{-0.12in}
\paragraph{Baselines.}
We benchmark against nine baseline methods. The first three are classic tabular modeling techniques: (1) Logistic regression (LogReg), (2) XGBoost~\cite{chen2016xgboost}, and (3) RandomForest~\cite{ho1995random}. Two subsequent baselines leverage unlabeled data: (4) SCARF~\cite{bahri2022scarf} and (5) STUNT~\cite{nam2023stunt}; however, practical acquisition of large amounts of unlabeled tabular data can be a significant bottleneck. Another modern approach, (6) TabPFN~\cite{hollmann2022tabpfn}, uses synthetic data generation spanning varied distributions for model pre-training. The last set of baselines applies LLM methodologies: (7) In-context learning (In-context)~\cite{wei2022emergent}, (8) TABLET~\cite{slack2023tablet}, and (9) TabLLM~\cite{hegselmann2023tabllm}. In-context learning operates by placing a limited number of demonstration samples in the prompt without any modification of LLM weights. TABLET supplements the prompts with externally produced hints derived from rule-based or prototype-based classifiers, in an effort to enrich the inference procedure. TabLLM fine-tunes LLM parameters efficiently via methods like IA3~\cite{liu2022few}, leveraging only a handful of labeled samples.

\vspace*{-0.12in}
\paragraph{Implementation details.}
\textsf{FeatLLM} adopts GPT-3.5 as its primary LLM backbone but maintains flexibility to function with various LLM architectures (refer to Appendix~\ref{sec:backbones} for experiments on alternate backbones). LLM inference operates with a temperature of 0.5 and the top-p parameter set at the default API value of 1. For ensemble construction and rule extraction, we use 20 and 10 as respective settings. The influence of these hyper-parameters is depicted in Figure~\ref{fig:hyperparameter2} and detailed further in Appendix~\ref{sec:hyper-parameter}. For the linear model, training employs the Adam optimizer at a learning rate of 0.01, running over 200 epochs, with optimal epochs selected via $k$-fold cross-validation. For baseline reproduction, GPT-3.5 serves as the model in in-context learning, while T0~\cite{sanh2022multitask} is used for TabLLM to align with its original design. TabLLM's protocol limits LLM usage to those with publicly released checkpoints, making GPT-3.5 unavailable within this benchmark. Conventional machine learning baselines, including LogReg, XGBoost, and RandomForest, have hyper-parameters tuned through grid search integrated with $k$-fold cross-validation. The chosen $k$ is set to 2 or 4 to ensure the presence of every class in the training partitions. Further implementation specifics are provided in Appendix~\ref{sec:baselines}.

\vspace*{-0.12in}
\paragraph{Results.}
Tables~\ref{tab:main1} and \ref{tab:main2} display comparative results across multiple datasets. Each experiment was conducted three times with distinct random splits for training and testing, and we report the mean and standard deviation of the resulting AUC metrics. \textsf{FeatLLM} either consistently outperforms all baselines or ranks second in most cases. Notably, our framework delivers nearly equivalent results to standard tabular methods employing 32-shot samples, while itself requiring only 4-shot samples (refer to Fig.~\ref{fig:compares}). While STUNT and TabPFN displayed marked improvement on specific datasets, their overall gains relative to established machine learning baselines were modest. Furthermore, LLM-based models including in-context learning, TABLET, and TabLLM, could not process tabular datasets featuring high dimensionality. By leveraging both bagging and ensembling, our framework demonstrates strong and reliable performance even in settings with a large number of features. Importantly, the bagging and ensemble strategies we introduce could be extended to other LLM-based techniques, but this extension would effectively double per-sample inference requirements, substantially increasing computational demand and rendering it less feasible in practice.

\subsection{Analysis \& Discussion}
\label{sec:analysis}
\paragraph{Ablation study.} We conducted ablation studies to evaluate the impact of the principal components on performance, considering the following scenarios: (1) \textsf{-Tuning}: eliminating the linear model's weight tuning, instead aggregating the binary features to produce the logit; (2) \textsf{-Ensemble}: removing the ensemble stage; (3) \textsf{-Description}: excluding feature description inputs; and (4) \textsf{-Reasoning}: not including the Step-1 procedure within the reasoning instruction when rules are generated.

Table~\ref{tab:ablation} illustrates how AUC is affected by each ablation, averaged over all datasets. Removing any of these components led to observable declines in model performance. Notably, the advantage conferred by weight tuning increases as the number of shots grows, signifying that, with more samples, the precise estimation of rule importance—and thereby tuning—becomes more beneficial. In contrast, incorporating feature descriptions and structured reasoning is most advantageous under few-shot regimes, implying that capitalizing on the prior knowledge embedded in LLMs is particularly valuable when labeled data is scarce. Across every scenario, the ensemble mechanism yields consistent performance improvements.

\vspace*{-0.12in}
\paragraph{Training \& inference time comparison.} We evaluate and contrast both training and inference times across different approaches. Our analysis utilizes the Adult dataset and a training subset comprising 16 instances. All models are executed on a single A100 GPU by default, with the exception of TabLLM, which utilizes four A100 GPUs for model parallelism. The timing metrics can be found in Table~\ref{tab:cost}. Remarkably, \textsf{FeatLLM} achieves low inference latency, on par with traditional baselines. On the other hand, Large Language Model-based techniques such as In-context learning, TABLET, and TabLLM incur substantially longer inference times. Regarding the training phase, \textsf{FeatLLM} aligns closely with its LLM-driven counterparts in duration—requiring just 30 API calls, a manageable computational cost that is also amenable to parallel execution.

\vspace*{-0.12in}
\paragraph{Quality of features generated by \textsf{FeatLLM}.}
To assess the practical value of rule-based features created by \textsf{FeatLLM}, we run a targeted experiment that measures how informative these features are in real-world tasks. Specifically, for each dataset, we compute the mean AUROC for the newly crafted features with respect to the target labels and then calculate the fraction of features achieving an AUROC above 0.5. Results are reported in Table~\ref{tab:quality}. The findings indicate that a significant proportion of the generated rules have clear predictive value for the target variable, effectively supporting the downstream learning task (see the “Before tuning” column). Furthermore, as the linear model undergoes training, rules that fail to demonstrate data relevance—specifically, those with associated learned weights beneath a specified threshold—are systematically excluded from the model (see the “After tuning” column). The threshold is fixed at 0.1, selected based on the distribution of learned weights observed in the experiments.

\vspace*{-0.12in}
\paragraph{Effect of adding spurious correlations.} To evaluate the capacity of different approaches to eliminate noisy spurious correlations under conditions of limited supervision, we performed a targeted study. We selected the Heart dataset as a basis and artificially appended randomly chosen columns from the Adult dataset that are not relevant to the Heart prediction task. Model accuracy was then monitored as we incrementally increased the number of these extraneous features. To ensure consistency, we standardized the scenario to 8 labeled samples per method and omitted the use of any unlabeled data, resulting in the exclusion of algorithms such as SCARF and STUNT from this experiment.

Figure~\ref{fig:spurious} presents the trajectories of model performance as more spurious correlations are introduced. Notably, \textsf{FeatLLM} exhibited superior resilience compared to baseline models, experiencing the smallest drop in accuracy due to the presence of irrelevant features. This effect can be attributed to our framework’s incorporation of prior knowledge, which enforces explicit associations between the prediction task and selected features during the rule extraction stage. Such a strategy inhibits the creation of rules grounded in extraneous columns, thereby promoting robustness. Quantitatively, we found that rules built upon the randomly added columns constituted only 1-2\% of the complete rule set. It is worth mentioning, however, that \textsf{FeatLLM} can still fall prey to learning spurious relationships and erroneously infer causality if highly similar but irrelevant features are appended (see Appendix~\ref{sec:spurious-appendix} for further discussion).

\vspace*{-0.12in}
\paragraph{Hyper-parameter analysis}
\textsf{FeatLLM} incorporates two primary hyper-parameters: the ensemble size and the number of rules extracted per class within each LLM query. We systematically evaluated their respective influences on the model’s efficacy. Our findings indicate that increasing the ensemble size enhances predictive performance, but the improvement plateaus at approximately 20 ensembles (refer to Fig.~\ref{fig:appendix-hyperparameter1} in the Appendix). Regarding the number of extracted rules, although differences did not reach statistical significance, setting this parameter at 10 was deemed most suitable, balancing prompt size with model accuracy (see Fig.~\ref{fig:hyperparameter2}). We hypothesize that raising the number of rules expands the dimensionality of inputs and increases informational content, but in the low-shot context, this can exacerbate overfitting.

\section{Conclusion}
We introduce \textsf{FeatLLM}, a method that advances LLM-based few-shot learning for tabular data. Unlike previous approaches that directly use LLMs for instance-level prediction, \textsf{FeatLLM} leverages LLMs to generate rule-based features—mirroring the process of human feature engineering. Employing rule generation in tandem with an ensemble strategy using bagging, \textsf{FeatLLM} achieves high predictive performance while incurring minimal inference cost. It operates solely through API access, requires no model training, and scales independently of tabular feature space size. These characteristics make \textsf{FeatLLM} an effective and practical solution for real-world applications demanding data efficiency.

While the current design of \textsf{FeatLLM} focuses on scenarios with sparse labeled data, future work will target extending its feature-creation paradigm to settings with larger datasets, investigate a broader array of feature types beyond rule-based constructions, and further promote interpretability to facilitate generalization to more diverse task domains.

\section{Broader Impact}
\textsf{FeatLLM} is intended to capitalize on the background knowledge encapsulated by LLMs, surpassing the conventional constraints of supervised tabular learning—especially under stringent data conditions. Our work contributes toward increasing data efficiency and curbing overfitting in contexts where labeled data is scarce, by incorporating the informative priors available in LLMs. This approach holds particular promise in domains like finance and healthcare, where obtaining labeled data is labor-intensive and expensive, yet pretrained language models often already contain significant relevant domain knowledge. However, as a critical avenue for future research and broader adoption, it is imperative to systematically examine how societal biases or inaccuracies present in LLM prior knowledge can shape, or even override, predictions—highlighting the need for thorough evaluation and mitigation strategies surrounding such influences.

\end{document}